%!TEX program = xelatex
\documentclass[lang=en,11pt,a4paper,citestyle=authoryear]{elegantpaper}

\title{HW2 - STAT 709}
\author{Wells Guan}

\usepackage{array,url,stix}
\usepackage{tikz}
\usepackage{tikz-cd}
\usepackage{enumitem}
\usepackage{amsmath}
\usepackage{amssymb}
\usepackage{listings}
\usepackage{subfigure}

\lstset{
    frame=single,
    columns=fullflexible,
    keepspaces=true,
    showstringspaces=false,
    breaklines=true,
    breakatwhitespace=false,
    basicstyle=\ttfamily\footnotesize,
    numbers=left,
    numberstyle=\tiny,
    numbersep=6pt
}
\newcommand{\Z}{\mathbb{Z}}
\newcommand{\R}{\mathbb{R}}
\newcommand{\N}{\mathbb{N}}
\newcommand{\B}{\mathcal{B}}
\newcommand{\E}{\mathbb{E}}
\renewcommand{\P}{\mathbb{P}}
\newcommand{\norm}[1]{\left\lVert#1\right\rVert}

\begin{document}

\maketitle
\subsection*{Notation and conventions}

All random variables are defined on a probability space
$(\Omega,\mathcal F,\P)$, and $\E$ denotes expectation with respect to
$\P$. For $1\leq p<\infty$, write
\[
\norm{X}_p=\bigl(\E|X|^p\bigr)^{1/p},
\qquad
\norm{X}_\infty
=\inf\{M\geq0:\P(|X|\leq M)=1\}.
\]
If $F$ is a cdf, then $dF$ denotes the probability measure induced by $F$,
so integrals with respect to $dF$ are Lebesgue--Stieltjes integrals. We use
$s\vee a=\max\{s,a\}$, and $G\sim N(0,1)$ means that $G$ is a standard
normal random variable. Statements about equality or boundedness of random
variables are understood almost surely when appropriate.



\subsection*{Problem 1. Ex.35. (a),(b)}

Let $\{a_n\}_{n\geq1}$ be a sequence of positive numbers satisfying
\[
\sum_{n=1}^{\infty}a_n=1,
\]
and let $\{P_n\}_{n\geq1}$ be a sequence of probability measures on a common
measurable space. Define
\[
P=\sum_{n=1}^{\infty}a_nP_n.
\]
\begin{enumerate}[label=(\alph*),leftmargin=*]
    \item Show that $P$ is a probability measure.
    \item For a measure $\nu$, show that $P_n\ll\nu$ for all $n$ if and only
    if $P\ll\nu$. Moreover, when $P\ll\nu$ and $\nu$ is $\sigma$-finite,
    show that
    \[
    \frac{dP}{d\nu}
    =\sum_{n=1}^{\infty}a_n\frac{dP_n}{d\nu}.
    \]
\end{enumerate}

\vspace{0.5em}
\textbf{Sol.}\par
(a) We know that $P$ is a measure from the Ex 1.7. in HW1. Now only need to check that
\[
P(\Omega) = \sum\limits_{n=1}^{\infty}a_nP_n(\Omega) = \sum\limits_{n=1}^{\infty} a_n = 1
\]
and we are done.\par
(b) To see the sufficiency, we know for any measurable set $A$,
\[P(A) \geq a_nP_n(A) \implies P_n(A )\leq P(A)/a_n\]
and hence if $\nu(A) = 0$, then $P_n(A) \leq P(A)/a_n = 0$ for any $n$, which means $P_n \ll \nu$ for any $n$. The necessity is easy to check by noticing if $\nu(A) = 0$, then $P_n(A) = 0$ for any $n$ and hence
\[
P(A) = \sum\limits_{n=1}^{\infty} a_nP_n(A) = \sum\limits_{n=1}^{\infty} a_n\cdot 0 = 0.
\]
To sum up, we showed the required equivalence.\par
If $P \ll \nu$, with $\nu$ $\sigma$-finite, we know that $P_n \ll \nu$ by the conclusion above. Denote
\[f = \sum\limits_{n=1}^{\infty} a_n \dfrac{dP_n}{d\nu}\]
and for any measurable set $A$, by the MCT, we know
\[
\int_A f d\nu = \sum\limits_{n=1}^{\infty} a_n \int_A \dfrac{dP_n}{d\nu} d\nu = \sum\limits_{n=1}^{\infty} a_nP_n(A) = P(A)
\]
and hence $f = \dfrac{dP}{d\nu}$ by the Radon-Nikodym Theorem.
\vspace{0.5em}

\subsection*{Problem 2. Ex.40.}

Let $X$ be a random variable having a continuous c.d.f. $F$. Show that
$Y = F(X)$ has the uniform distribution $U(0, 1)$.\par
\vspace{0.5em}
\textbf{Sol.}\par
Define
\[
F^{-1}(a) = \sup\{y, F^-(y) < a\}
\]
and we know $F(F^{-1}(a)) \geq a$ by the right-continuity. Notice
\[
P(Y \geq a) =  P(F(X) \geq a) = P(X \geq F^{-1}(a)) = 1 - F^-(F^{-1}(a))
\]
where 
\[
F^{-}(F^{-1}(a)) \leq a
\]
since $F^-$ is left-continuous. However, $F^- = F$ since $F$ continuous and hence $F^{-}(F^{-1}(a)) = a$. Then we know
\[
P(Y\geq a) = 1-a
\]
and we are done.
\vspace{0.5em}

\subsection*{Problem 3. Ex.55.}
Let $X$ be a random variable. Show that
\begin{enumerate}[label=(\arabic*),leftmargin=*]
    \item[(a)] if $EX$ exists, then
      \[
      EX = \int_0^{\infty} P(X > x)dx - \int_{-\infty}^0 P(X\leq x) dx
      \]
    \item[(b)] if $X$ has range $\N$, then \[EX = \sum\limits_{n\geq 1} P(X\geq n)\]
\end{enumerate}
\vspace{0.5em}
\textbf{Sol.}\par
(a) Notice that by the Fubini's Theorem, we have
\[
\begin{aligned}
  EX &= \int_{-\infty}^{\infty} yd\mu(y) \\
  &= \int_{(0,\infty)} yd\mu(y) - \int_{(-\infty,0]} (-y)d\mu(y)\\
  &= \int_{(0,\infty)} \int_{(0,y)} dx d\mu(y) - \int_{(-\infty,0]} \int_{[y,0)} dx d\mu(y) \\
  &= \int_0^{\infty} \int_{(x,\infty)} d\mu(y)dx - \int_{-\infty}^0 \int_{(-\infty,x]} d\mu(y)dx     \\
  &= \int_0^{\infty} P(X>x) dx - \int^0_{-\infty}P(X\leq x)dx
\end{aligned}
\]
where $\mu$ is the induced measure by $X$.\par
(b) Use the inequality in (a) directly and we have
\[
\begin{aligned}
  EX &= \int_0^{\infty} P(X>x) dx - \int^0_{-\infty}P(X\leq x)dx \\
  &= \sum\limits_{n=0}^{\infty} \int_n^{n+1} P(X \geq n+1) - \int_0^0 P(X = 0) dx \\
  &= \sum\limits_{n\geq 1} P(X\geq n)
\end{aligned}
\]
and we are done.
\vspace{0.5em}

\subsection*{Problem 4. Ex.83.}
 Let $X$ be an integrable random variable with a Lebesgue p.d.f. $f_X$
and let $Y = g(X)$, where $g$ is a function with positive derivative on $(0,\infty)$ and $g(x) = g(-x)$. Find an expression for $E(X|Y)$ and verify
that it is indeed the conditional expectation.\par
\vspace{0.5em}
\textbf{Sol.}\par
Let
\[
h=\left(g\big|_{[0,\infty)}\right)^{-1}
\]
be the inverse of the restriction of $g$ to $[0,\infty)$. Since $g$ is even
and strictly increasing on $[0,\infty)$, we have
\[
h(Y)=h(g(X))=|X|
\qquad \text{a.s.}
\]

For $u\geq 0$, define
\[
r(u)=
\begin{cases}
\dfrac{f_X(u)-f_X(-u)}
      {f_X(u)+f_X(-u)},
& f_X(u)+f_X(-u)>0,\\[1.2ex]
0,
& f_X(u)+f_X(-u)=0.
\end{cases}
\]
and we prove that
\[
E(X\mid Y)=h(Y)r(h(Y)).
\]
Denote
\[
Z=h(Y)r(h(Y)).
\]
Since $Z$ is a Borel function of $Y$, it is $Y$-measurable. Moreover,
$|r(u)|\leq 1$, and hence
\[
|Z|\leq h(Y)=|X|.
\]
Since $X$ is integrable, $Z$ is also integrable.

Let $\varphi$ be any bounded Borel function. Using the fact that
$g(-u)=g(u)$, we obtain
\[
\begin{aligned}
E\bigl[\varphi(Y)X\bigr]
&=\int_{\mathbb R}\varphi(g(x))x f_X(x)\,dx\\
&=\int_0^\infty \varphi(g(u))u f_X(u)\,du
  -\int_0^\infty \varphi(g(u))u f_X(-u)\,du\\
&=\int_0^\infty
  \varphi(g(u))u\bigl(f_X(u)-f_X(-u)\bigr)\,du.
\end{aligned}
\]
On the other hand,
\[
\begin{aligned}
E\bigl[\varphi(Y)Z\bigr]
&=\int_{\mathbb R}
  \varphi(g(x))h(g(x))r(h(g(x)))f_X(x)\,dx\\
&=\int_0^\infty
  \varphi(g(u))u r(u)
  \bigl(f_X(u)+f_X(-u)\bigr)\,du\\
&=\int_0^\infty
  \varphi(g(u))u
  \bigl(f_X(u)-f_X(-u)\bigr)\,du.
\end{aligned}
\]
Therefore,
\[
E\bigl[\varphi(Y)Z\bigr]
=
E\bigl[\varphi(Y)X\bigr]
\]
for every bounded Borel function $\varphi$ and we are done.
\vspace{0.5em}


\subsection*{Problem 5. Integration by parts}

Let $a<b$ and $f,g\in C^1([a,b])$, the space of continuously differentiable
functions on $[a,b]$. The usual integration-by-parts formula is
\[
\int_a^b f'(x)g(x)\,dx
=f(b)g(b)-f(a)g(a)-\int_a^b f(x)g'(x)\,dx.
\]
We can regard $f'(x)\,dx$ and $g'(x)\,dx$ as defining signed measures on
$[a,b]$. If the derivatives are nonnegative, these are nonnegative measures.
We now generalize the formula to a probability measure with cdf $F$. Recall
that probability measures on $\R$ correspond one-to-one to cdfs. For
$g\in C^1([a,b])$, prove that
\[
\int_{(a,b]}g(x)\,dF(x)
=F(b)g(b)-F(a)g(a)-\int_a^b g'(x)F(x)\,dx.
\]

\vspace{0.5em}
\textbf{Sol.}\par
Notice the by the Fubini's theorem we have
\[
\begin{aligned}
&\int_a^b g'(x)F(x)\,dx+
  \int_{(a,b]}g(x)\,dF(x)\\
&=\int_a^b g'(x)\int_{(-\infty,x]}dF(s)\,dx
  +\int_{(a,b]}g(s)\,dF(s)\\
&=\int_{(-\infty,b]}\int_{s\vee a}^{b}g'(x)\,dx\,dF(s)
  +\int_{(a,b]}g(s)\,dF(s)\\
&=\int_{(-\infty,b]}\bigl(g(b)-g(s\vee a)\bigr)\,dF(s)
  +\int_{(a,b]}g(s)\,dF(s)\\
&=\int_{(-\infty,a]}\bigl(g(b)-g(a)\bigr)\,dF(s)
  +\int_{(a,b]}\bigl(g(b)-g(s)\bigr)\,dF(s)
  +\int_{(a,b]}g(s)\,dF(s)\\
&=\bigl(g(b)-g(a)\bigr)F(a)
  +g(b)\bigl(F(b)-F(a)\bigr)\\
&=g(b)F(b)-g(a)F(a),
\end{aligned}
\]
and we are done.

\vspace{0.5em}
\subsection*{Problem 6. Moment bounds}

Let $X$ be a real-valued random variable.

\begin{enumerate}[label=(\arabic*),leftmargin=*]
    \item Suppose $0<p<q$ and $\mathbb E|X|^q<\infty$. Use H\"older's
    inequality to show that $\mathbb E|X|^p<\infty$ and
    \[
    \bigl(\mathbb E|X|^p\bigr)^{1/p}
    \leq \bigl(\mathbb E|X|^q\bigr)^{1/q}.
    \]

    \item Suppose
    \[
    M=\sup_{p\geq1}\bigl(\mathbb E|X|^p\bigr)^{1/p}<\infty.
    \]
    Prove that $\mathbb P(|X|\leq M)=1$.

    \item[(3)] Let $G\sim N(0,1)$. Prove that there is a
    constant $c>0$ such that
    \[
    \lim_{n\to\infty}
    \frac{\bigl(\mathbb E|G|^n\bigr)^{1/n}}{\sqrt n}=c,
    \]
    and determine $c$.
\end{enumerate}

\vspace{0.5em}
\textbf{Sol.}\par

\textbf{(1)} By H\"older's inequality, with conjugate exponents
$q/p$ and $q/(q-p)$,
\[
\begin{aligned}
\mathbb E|X|^p
&=\mathbb E\bigl(|X|^p\cdot1\bigr)\\
&\leq
\left(\mathbb E|X|^{p(q/p)}\right)^{p/q}
\left(\mathbb E1^{q/(q-p)}\right)^{(q-p)/q}\\
&=\left(\mathbb E|X|^q\right)^{p/q}.
\end{aligned}
\]
Therefore $\mathbb E|X|^p<\infty$, and taking $p$th roots gives
\[
\left(\mathbb E|X|^p\right)^{1/p}
\leq\left(\mathbb E|X|^q\right)^{1/q}.
\]

\textbf{(2)} Fix $\varepsilon>0$ and write
\[
\alpha=\mathbb P(|X|\geq M+\varepsilon).
\]
If $\alpha>0$, then for every $p\geq1$,
\[
\begin{aligned}
\mathbb E|X|^p&\geq(M+\varepsilon)^p\alpha,\\
\implies \qquad
\left(\mathbb E|X|^p\right)^{1/p}
&\geq(M+\varepsilon)\alpha^{1/p}.
\end{aligned}
\]
Let $p\to\infty$ contradicts the definition of $M$, and hence
$\mathbb P(|X|\geq M+\varepsilon)=0$ for every $\varepsilon>0$, and therefore
\[
\mathbb P(|X|>M)=0,
\qquad
\mathbb P(|X|\leq M)=1.
\]

\textbf{(3)} Let $m_n=\mathbb E|G|^n$. Integration by parts gives, for
$n\geq2$,
\[
\begin{aligned}
m_n
&=\frac1{\sqrt{2\pi}}
  \int_{\R}|x|^n e^{-x^2/2}\,dx\\
&=\frac1{\sqrt{2\pi}}
  \int_{\R}|x|^{n-2}x\bigl(xe^{-x^2/2}\bigr)\,dx\\
&=(n-1)\frac1{\sqrt{2\pi}}
  \int_{\R}|x|^{n-2}e^{-x^2/2}\,dx\\
&=(n-1)m_{n-2}.
\end{aligned}
\]
In particular,
\[
m_{2k}=(2k-1)(2k-3)\cdots1
=\frac{(2k)!}{2^k k!}.
\]
By Stirling's formula,
\[
\begin{aligned}
\frac{m_{2k}^{1/(2k)}}{\sqrt{2k}}
=\frac1{\sqrt{2k}}
 \left(\frac{(2k)!}{2^k k!}\right)^{1/(2k)} &\simeq  \dfrac{1}{\sqrt{2k}}\left(\dfrac{(2k)!}{\left(\sqrt{2\pi(2k)}\left(\dfrac{2k}{e}\right)\right)^{2k}}\dfrac{2^k\left(\sqrt{2\pi(k)}\left(\dfrac{k}{e}\right)\right)^{k}}{2^kk!}\right)^{1/2k}\dfrac{\sqrt{2}^{1/2k}\sqrt{2k}}{\sqrt{e}}\\
& \longrightarrow \frac1{\sqrt e}.
\end{aligned}
\]
The $L^p$ norms are increasing, so
\[
m_{2k}^{1/(2k)}
\leq m_{2k+1}^{1/(2k+1)}
\leq m_{2k+2}^{1/(2k+2)}.
\]
and since \[\lim\limits_{k\to\infty} \dfrac{\sqrt{2k+1}}{\sqrt{2k}} = \lim\limits_{k\to\infty} \dfrac{\sqrt{2k+1}}{\sqrt{2k+2}} =1,\] we know both bounds have the same limit
$e^{-1/2}$. Therefore, we have
\[
\lim_{n\to\infty}
\frac{\bigl(\mathbb E|G|^n\bigr)^{1/n}}{\sqrt n}
= c,
\]
with $c=e^{-1/2}$.

\end{document}
